% !TeX root = ../papers/nc-p3.tex
\section{周期一の共通窓における全成分の局所方程式}
\label{sec:torus-component-crossing}
\begin{theorem}\label{thm:torus-component-crossing}
単重み曲面と単重み曲線の共通点 $w_*$、すなわち $(b,c)=(1/2,1)$ において、
トーラスを仮定しない二切断階数6軌跡の軌道横断完備局所環は
\[
 k[[u,v,z,t_1,\ldots,t_{14}]]/(uv,uz)
\]
である。$u,v,z$ は重み0、残りは正負各側で $1,2,3,3,4,4,5$ の絶対重みを持つ。
アフィンテンソルの完備局所環にはさらに60個の滑らかな軌道変数を加える。
従ってこの点を通る全体成分はちょうど二つで、双方ともここで滑らか、交叉も滑らかである。
二成分のアフィン次元は76と75、共通部分は74であり、交叉のスキームに重複はない。
この閉軌道に特殊化する窓は、閉軌道自身のほか、一つの $\mathbb P(1,2,3,3,4,4,5)$ で分類される。
これらは全て有限点グラフの標準正則窓で両成分の共通部分にある。
\end{theorem}
\begin{proof}[重み0におけるスキーム構造]
前節の三十八単元の規格化を使う。重み0の基底変換に横断な十二主元を1にし、
重み3ブロックのSchur方程式から $d,h,f,x,y$ を解いた残りの三つの単元 $a,b,c$ を取る。
全ての消去の分母は $w_*$ で単元である。新しいエタール座標を
\[
 u=ab-c,\qquad v=a(1-b)-c,\qquad z=c-1
\]
とする。逆に $c=1+z$, $a=u+v+2c$, $b=(u+c)/a$。
重み0のイデアルは $I_0=(uv,uz)$ である。実際一式は単元倍の $uv$ となり、
別の式の $v=0$ での制限は $uz$ の単元倍となる。全式は $u=0$ 上と $v=z=0$ 上でそれぞれ恒等零である。
よって包含 $I_0\subset (u)\cap(v,z)=(uv,uz)$ を得る。
逆向きには後者の式を $u$ で割って $v$ を法とした剰余が $z$ の単元倍であることと $uv\in I_0$ を使う。
\end{proof}

\begin{proof}[全横断環での共通因子]
前節と同じ一次条件で接空間の軌道商は17次元、うち重み0が3次元で正負は上表のままである。
両極の局所環上の四文字積の階数は19である。これらと上記単元の検査は
\path{computations/verify_torus_component_crossing.py} による。
有限平坦な点環の19恒等式により、軌道を除き線形独立な微分を持つ方程式を消去した後では、
17変数正則完備局所環 $L$ に残るイデアル $I$ は高々二本で生成される。
実際アフィンでは独立微分が $256-60-17=179$、局所生成元が高々 $200-19=181$ 本である。
トーラスは線形還元的なので、この消去では独立な微分の同変な持上げを選べる。
固定部分の軌道は対角基底変更の軌道であり、その完備横断環は上の十二主元による規格化と同型である。
従って $L$ の非零重み変数を零にしたときの $I$ の像は $I_0$ である。
$I_0$ に二生成元が必要なので、$I/\mathfrak m_L I$ は二次元で $I_0/\mathfrak m I_0$ へ同型に写る。
特に二つの生成元 $f_1,f_2$ を重み0に選べる。

前節の76次元成分がこの点を通るため、$I$ は $L$ の高さ1の素イデアルに入る。
正則局所環は一意分解環なので $f_1,f_2$ に非単元共通因子 $F$ がある。
重み0の三変数へのその制限は、零でも単元でもなく、$uv,uz$ の最大公約元 $u$ を割る。
従って制限は $u$ の単元倍であり、$F$ は滑らかな助変数である。
別の共通素因子もその二乗もあり得ない。この一意な主イデアルはトーラス安定であり、
$F$ の重み0部分も同じイデアルを生成するので $F$ を不変にできる。
$f_i=F g_i$ と置けば $g_i$ の固定部分での制限は $(v,z)$ を生成する。
従って $F,g_1,g_2$ は独立な一次部を持つ三つの重み0の変数に取れ、残りの変数を
同変に補えば主張した完備局所方程式になる。既に相異なる二つの実際の成分がここを通るため、
二つの完備既約枝はその二成分に一つずつ対応する。
\end{proof}

\begin{proof}[特異点の零ファイバー]
スライスを滑らかな線形表現に局所閉埋込みし、独立な微分を持つ方程式の持上げを
安定群同変に取ると、その零点はスライスを含む滑らかな17次元の周囲を与える。
ここで安定群はトーラスとそれを反転する有限成分からなり線形還元的である。
周囲のスライスの零ファイバーは前節の強エタール性から正・負の各七次元空間の和集合になる。
この全体が実際のテンソルスライスに入る。実際、残るイデアルの極小生成空間は上で示した通りトーラス重み0だけなので、
アフィン近傍でトーラス不変な生成元に持ち上げられる。これらは原点で消え、原点を極限に持つ全点で消える。
同じ議論を各枝のイデアルに用いることもできる。上の完備表示で両枝の極小生成元は重み0であり、
完備化前の各枝の極小生成空間も同じである。従って両枝とも零ファイバー全体を含む。
向きを保つ安定群は対角トーラスで、重み反転が正負を交換するため、一つの重み付き射影空間を得る。
\end{proof}
